\par\noindent\textbf{(1)} 单镜振幅反射系数r且|r|²=R，透射幅t满足|t|²=1−R。相邻透射光束多一次往返相因子 R e\textasciicircum{}\{2iδ\}，总幅 E\_t/E\_in=t²e\textasciicircum{}\{iδ\}Σ\_\{j≥0\}(Re\textasciicircum{}\{2iδ\})\textasciicircum{}j=(1−R)e\textasciicircum{}\{iδ\}/(1−Re\textasciicircum{}\{2iδ\})，须相干求和。
\par\noindent\textbf{(2)} 取幅模平方得 T(δ)=(1−R)²/[(1−R)²+4Rsin²δ]=1/[1+F sin²δ]，F=4R/(1−R)²。透射最大为1，条件 δ=mπ 即光学往返相位为2mπ；无损两镜的完全峰透射与R多大无关。
\par\noindent\textbf{(3)} 相邻纵模频率间隔FSR=c/(2nL)。高精细度F≫1时共振邻域 sin²δ≈(δ−mπ)²，半高时|δ−mπ|=1/√F；换算全宽 Δν≈FSR·2/(π√F)，精细度𝓕=FSR/Δν≈π√R/(1−R)。
\par\noindent\textbf{(4)} 光子在腔内平均多次往返，驻留时间尺度与 1/(1−R) 成正比，共振透射的群延迟与频率导数有关，不能简单把腔内能量除输入功率定义为任意频率群时延。R→0恢复单次透射、Airy无选择性；R→1峰变窄但耦合带宽趋零。
